% GENERATED FILE. Edit canonical lessons, then run npm run generate:books.
% canonical-source-hash: fnv1a64:cfc24aafd4a20efc
% canonical-lessons: SA-C150-bhesajam, SA-C150-vranah, SA-C150-sothah, SA-C150-vamanam, SA-C150-udarapida

\chapter{Remedy, Wound, Swelling, Vomiting, Stomach ache}
\label{ch:sa-remedy-wound-swelling-vomiting-stomach-ache}
\hlchaptermodality{150}

\begin{chapteropening}
\textbf{By the end of this chapter:} \emph{I can say a remedy, a wound, a swelling, vomiting and a stomach ache.}

Complete the last lesson of chapter 150: I can say a remedy, a wound, a swelling, vomiting and a stomach ache.
\end{chapteropening}

\section[\texorpdfstring{\sk{भेषजम्} (bheṣajam)}{भेषजम् (bheṣajam)}]{\sk{भेषजम्} (bheṣajam) — a remedy}

\label{lesson:SA-C150-bhesajam}



\begin{quote}
Before the new one: say the Sanskrit for a clerk, then the Sanskrit for a leader.
\end{quote}

\subsection*{You'll want to know: \sk{भेषजम्}}
\textbf{\sk{भेषजम्}} — \emph{bheṣajam} — \textquotedblleft{}a remedy\textquotedblright{}.

\begin{grammarlens}[title={What you've built}]
One more word for this chapter.
\end{grammarlens}

\subsection*{Your turn}
\begin{itemize}
\raggedright
  \item \emph{Say it:} \emph{bheṣajam}
  \item \emph{Say it:} \emph{bheṣajam}, once more
  \item \emph{Say it:} say \emph{bheṣajam}
  \item \emph{Recall:} say the Sanskrit for a clerk, then the Sanskrit for a leader, then say \emph{bheṣajam} again
\end{itemize}

\begin{tcolorbox}[breakable,colback=teal!4,colframe=teal!35!black,title={Before you move on}]
What does \sk{भेषजम्} mean? (\textquotedblleft{}A remedy\textquotedblright{}.) Say it once more.
\end{tcolorbox}
\section[\texorpdfstring{\sk{व्रणः} (vraṇaḥ)}{व्रणः (vraṇaḥ)}]{\sk{व्रणः} (vraṇaḥ) — a wound}

\label{lesson:SA-C150-vranah}



\begin{quote}
Before the new one: say the Sanskrit for a leader, then the Sanskrit for a remedy.
\end{quote}

\subsection*{You'll want to know: \sk{व्रणः}}
\textbf{\sk{व्रणः}} — \emph{vraṇaḥ} — \textquotedblleft{}a wound\textquotedblright{}.

\begin{grammarlens}[title={What you've built}]
One more word for this chapter.
\end{grammarlens}

\subsection*{Your turn}
\begin{itemize}
\raggedright
  \item \emph{Say it:} \emph{vraṇaḥ}
  \item \emph{Say it:} \emph{vraṇaḥ}, once more
  \item \emph{Say it:} say \emph{vraṇaḥ}
  \item \emph{Recall:} say the Sanskrit for a leader, then the Sanskrit for a remedy, then say \emph{vraṇaḥ} again
\end{itemize}

\begin{tcolorbox}[breakable,colback=teal!4,colframe=teal!35!black,title={Before you move on}]
What does \sk{व्रणः} mean? (\textquotedblleft{}A wound\textquotedblright{}.) Say it once more.
\end{tcolorbox}
\section[\texorpdfstring{\sk{शोथः} (śothaḥ)}{शोथः (śothaḥ)}]{\sk{शोथः} (śothaḥ) — a swelling}

\label{lesson:SA-C150-sothah}



\begin{quote}
Before the new one: say the Sanskrit for a remedy, then the Sanskrit for a wound.
\end{quote}

\subsection*{You'll want to know: \sk{शोथः}}
\textbf{\sk{शोथः}} — \emph{śothaḥ} — \textquotedblleft{}a swelling\textquotedblright{}.

\begin{grammarlens}[title={What you've built}]
One more word for this chapter.
\end{grammarlens}

\subsection*{Your turn}
\begin{itemize}
\raggedright
  \item \emph{Say it:} \emph{śothaḥ}
  \item \emph{Say it:} \emph{śothaḥ}, once more
  \item \emph{Say it:} say \emph{śothaḥ}
  \item \emph{Recall:} say the Sanskrit for a remedy, then the Sanskrit for a wound, then say \emph{śothaḥ} again
\end{itemize}

\begin{tcolorbox}[breakable,colback=teal!4,colframe=teal!35!black,title={Before you move on}]
What does \sk{शोथः} mean? (\textquotedblleft{}A swelling\textquotedblright{}.) Say it once more.
\end{tcolorbox}
\section[\texorpdfstring{\sk{वमनम्} (vamanam)}{वमनम् (vamanam)}]{\sk{वमनम्} (vamanam) — vomiting}

\label{lesson:SA-C150-vamanam}



\begin{quote}
Before the new one: say the Sanskrit for a wound, then the Sanskrit for a swelling.
\end{quote}

\subsection*{You'll want to know: \sk{वमनम्}}
\textbf{\sk{वमनम्}} — \emph{vamanam} — \textquotedblleft{}vomiting\textquotedblright{}.

\begin{grammarlens}[title={What you've built}]
One more word for this chapter.
\end{grammarlens}

\subsection*{Your turn}
\begin{itemize}
\raggedright
  \item \emph{Say it:} \emph{vamanam}
  \item \emph{Say it:} \emph{vamanam}, once more
  \item \emph{Say it:} say \emph{vamanam}
  \item \emph{Recall:} say the Sanskrit for a wound, then the Sanskrit for a swelling, then say \emph{vamanam} again
\end{itemize}

\begin{tcolorbox}[breakable,colback=teal!4,colframe=teal!35!black,title={Before you move on}]
What does \sk{वमनम्} mean? (\textquotedblleft{}Vomiting\textquotedblright{}.) Say it once more.
\end{tcolorbox}
\section[\texorpdfstring{\sk{उदरपीडा} (udarapīḍā)}{उदरपीडा (udarapīḍā)}]{\sk{उदरपीडा} (udarapīḍā) — a stomach ache}

\label{lesson:SA-C150-udarapida}



\begin{quote}
Before the new one: say the Sanskrit for a swelling, then the Sanskrit for vomiting.
\end{quote}

\subsection*{You'll want to know: \sk{उदरपीडा}}
\textbf{\sk{उदरपीडा}} — \emph{udarapīḍā} — \textquotedblleft{}a stomach ache\textquotedblright{}.

\begin{grammarlens}[title={What you've built}]
One more word for this chapter.
\end{grammarlens}

\subsection*{Your turn}
\begin{itemize}
\raggedright
  \item \emph{Say it:} \emph{udarapīḍā}
  \item \emph{Say it:} \emph{udarapīḍā}, once more
  \item \emph{Say it:} say \emph{udarapīḍā}
  \item \emph{Recall:} say the Sanskrit for a swelling, then the Sanskrit for vomiting, then say \emph{udarapīḍā} again
\end{itemize}

\begin{tcolorbox}[breakable,colback=teal!4,colframe=teal!35!black,title={Before you move on}]
What does \sk{उदरपीडा} mean? (\textquotedblleft{}A stomach ache\textquotedblright{}.) Say it once more.
\end{tcolorbox}
